   \batchmode
   \documentstyle[12pt]{article}
   \addtolength{\textwidth}{+1.05truein}
   \addtolength{\textheight}{+1.75truein}
   \addtolength{\hoffset}{-0.5truein}
   \addtolength{\voffset}{-0.85truein}
   
   \pagestyle{empty}
   
   \begin{document}
   
   \title{Semiconductor Wafer Manufacturing \\ Data Format Specification}
   
   \author{Gerry Feigin\thanks{With the help of: Sarah Hood, Frank Chance, and
   Dan Connors.}\\
   {\em IBM Research Division} \\
   {\em T. J. Watson Research Center} \\
   {\em P. O.  Box 218}   \\
   {\em Yorktown Heights, NY  10598}\\
   \\
   {John Fowler} \\
   {Jennifer Robinson} \\
   {\em SEMATECH} \\
   {\em 2706 Montopolis Drive} \\
   {\em Austin, TX 78741} \\
   \\
   {Robert Leachman} \\
   {\em University of California} \\
   {\em Dept. of Industrial Engineering and Operations Research} \\
   {\em Berkeley, CA 94720}  \\
   \\
   }
   \date{July 19, 1994}
   \maketitle
   
   %\begin{abstract}
   %\end{abstract}
   
   
   %\newpage
   %\setcounter{page}{2}
   \pagestyle{plain}
   
   \baselineskip 17pt plus 2pt minus 2pt
   
   \section{Introduction}
   This document describes the format for the semiconductor manufacturing
   line data as well as relevant comments and
   information about the data.  For each data set, there are 5 data files
   (or tables) as well as a sixth comments file.
   The name of each file for a given data set
   should have a prefix common to all the files for this data set
   followed by a unique suffix to identify the file.   For example,
   if the data set were \verb+ibm+, then the the files corresponding to this 
   data set would have the names \verb+ibm.<suffix>+.
   \\
   \\
   Below is a list of the
   six files, along with their suffixes and a brief description
   of the data they contain.
   
   \begin{center}
   \begin{tabular}{|l|c|l|}  \hline
   {\em File} & {\em Suffix ID} & {\em Description} \\ \hline
   Process Route &  pr &    Process route information for
                           all processes \\
   Rework Sequences & rw &   Information on rework sequences \\
   Tool Set   &     ts &   Information on tools  \\
   Operator Set  &  os  &  Information on operators  \\
   Volume/release & vr &   Release rate information \\ 
   Comment File  &  cf  &  General comments and sample run 
                          results  \\  \hline 
   \end{tabular}
   \end{center}
   
   \section{File descriptions}
   In this section, a detailed description is given
   of the format for each of the files described in the previous
   section. For each of the 5 data files, a list of all the fields for
   each record is given, and for each field, its type (e.g. integer,
   float, etc.), its width, and a brief description of its contents
   are specified.
   A description of the comments file is also provided. 

   
   \subsection{Process Route File}
   
   \begin{center}
   \begin{tabular}{|l|c|c|l|}  \hline
   {\em Field Name} &  {\em Type} & {\em Width } & {\em Remarks} \\ \hline
   Process Flow ID  & String & 10 & Unique process flow ID \\
   Step ID & String & 10 & Process step ID \\
   Operation Description  & String & 25 & A brief operation description \\
   Tool Set ID & String & 10 & Assume
   1 tool needed from this tool group \\
    & & & for this operation \\
   Operator Set ID & String & 15 &  Assume 1 operator needed for this
   operation \\
    & & & (see operator set file for more info) \\
   Load Time & Float & 10 & Time to load wafers/lot/batch into tool \\
   Unload Time & Float & 10 & Time to unload wafers/lot/batch from tool \\
   Time per Wafer in Process & Float & 10 & Processing time per wafer (excluding 
   travel \\
    & & & time within tool) (as appropriate) \\
   Wafer Travel Time & Float & 10 & Travel time within tool (as appropriate)\\
   Time per Lot & Float & 10 & Lot processing time (as appropriate) \\
   Time per Batch & Float & 10 & Batch processing time (as appropriate) \\
   Min Batch Size & Integer & 5 & Minimum batch size (in wafers) \\
   Max Batch Size & Integer & 5 & Maximum batch size (in wafers) \\
   Batch ID & String & 10 & Used to identify which operations \\
    & & & can be batched together \\
   Time per Spec. Setup & Float & 10 & Setup time required for changing from one 
   \\
    & & &  process flow step (recipe spec.) to another \\
   Time per Group Setup & Float & 10 & Setup time required for changing
   from \\
    & & & one ``group'' to another where group is specified \\
    & & & by the Setup Group ID  \\
   Setup Group ID & String & 15 & See above \\
   Lot Scrap Probability & Float & 10 & Prob. an entire lot is scrapped \\
    & & & after this operation \\
   Wafer Scrap Probability & Float & 10 & Prob. a wafer in a lot is
   scrapped \\
    & & & after this operation \\
   Lot Rework Probability & Float & 10 & Prob. an entire lot is
   sent for rework after \\
    & & & this operation given it has not been scrapped \\
   Wafer Rework Probability & Float & 10 & Prob. a wafer in a lot
   is sent for rework given it \\
    & & & is has not been scrapped \\
   Rework Sequence ID & String & 10 & Rework sequence to follow \\
   Rework Return Step ID & String & 10 & Step to which wafers return
   after rework \\
   Travel Time & Float & 10 & Time to travel to next operation process step \\
   Travel Time Operator ID & String & 15 &  Operator set
   needed for travel \\ \hline
   \end{tabular}
   \end{center}

   \newpage
   
   \subsection{Rework Sequences File}
   The format for a rework sequence record is almost identical
   to the corresponding format given for the process route.  The
   difference  is that instead of a Process ID, there is
   a rework sequence ID.
   
   \begin{center}
   \begin{tabular}{|l|c|c|l|}  \hline
   {\em Field Name} &  {\em Type} & {\em Width } & {\em Remarks} \\ \hline
   Rework Sequence ID  & String & 10 & Unique product family ID \\
   Step ID & String & 10 & Process step ID \\
   Operation Description  & String & 25 & A brief operation description \\
   Tool Set ID & String & 10 & Assume
   1 tool needed from this tool set \\
    & & & for this operation \\
   Operator Set ID & String & 15 &  Assume 1 operator needed for this
   operation \\
    & & & (see operator set file for more info) \\
   Load Time & Float & 10 & Time to load wafers/lot/batch into tool \\
   Unload Time & Float & 10 & Time to unload wafers/lot/batch from tool \\
   Time per Wafer in Process & Float & 10 & Processing time per wafer (excluding 
   travel \\
    & & & time within tool) (as appropriate) \\
   Wafer Travel Time & Float & 10 & Travel time within tool (as appropriate)\\
   Time per Lot & Float & 10 & Lot processing time (as appropriate) \\
   Time per Batch & Float & 10 & Batch processing time (as appropriate) \\
   Min Batch Size & Integer & 5 & Minimum batch size (in wafers) \\
   Max Batch Size & Integer & 5 & Maximum batch size (in wafers) \\
   Batch ID & String & 10 & Used to identify which operations \\
    & & & can be batched together \\
   Time per Spec. Setup & Float & 10 & Setup time required for changing from one 
   \\
    & & &  process flow step (recipe spec.) to another \\
   Time per Group Setup & Float & 10 & Setup time required for changing
   from \\
    & & & one ``group'' to another where group is specified \\
    & & & by the Setup Group ID  \\
   Setup Group ID & String & 15 & See above \\
   Lot Scrap Probability & Float & 10 & Prob. an entire lot is scrapped \\
    & & & after this operation \\
   Wafer Scrap Probability & Float & 10 & Prob. a wafer in a lot is
   scrapped \\
    & & & after this operation \\
   Lot Rework Probability & Float & 10 & Prob. an entire lot is
   sent for rework after \\
    & & & this operation given it has not been scrapped \\
   Wafer Rework Probability & Float & 10 & Prob. a wafer in a lot
   is sent for rework given it \\
    & & & is has not been scrapped \\
   Rework Sequence ID & String & 10 & Rework sequence to follow \\
   Rework Return Step ID & String & 10 & Step to which wafers return
   after rework \\
   Travel Time & Float & 10 & Time to travel to next operation process step \\
   Travel Time Operator ID & String & 15 &  Operator set
   needed for travel \\ \hline
   \end{tabular}
   \end{center}
   
   \subsection{Tool Set File}
   
   \begin{center}
   \begin{tabular}{|l|c|c|l|}  \hline
   {\em Field Name} &  {\em Type} & {\em Width } & {\em Remarks} \\ \hline
   Tool Set ID & String & 10 & Tool set identifier \\
   Tool Description & String & 25 & Name or description of tool \\
   Quantity & Integer & 5 & Number of (identical) tools in tool set \\
   Operator Load Fraction & Float & 10 & Fraction of time operator is needed \\
    & & & for lot loading \\
   Operator Unload Fraction & Float & 10 & Fraction of time operator is needed \\
    & & & for lot unloading \\
   Operator Process Fraction & Float & 10 & Fraction of time operator is
   needed \\
    & & & for lot processing \\
   Down Time \#1 Description & String & 25 & Description of tool down time
   type 1 \\
   Down Time \#1 Type & Integer & 5 & 0 = time based; 1 = run based \\
   Time or Runs Between \#1 & Float & 15 & Mean time (runs) between this \\
    & & & down time event \\
   Duration \#1 & Float & 10 & Duration of this down time event \\
   Down Time \#1 Operator Set ID & String & 15 & Operator needed from
    this set\\ 
    & & & during this down time event\\
   Down Time \#2 Description & String & 25 & Description of tool down time
   type 2 \\
   Down Time \#2 Type & Integer & 5 & 0 = time based; 1 = run based \\
   Time or Runs Between \#2 & Float & 15 & Mean time (runs) between this \\
    & & & down time event \\
   Duration \#2 & Float & 10 & Duration of this down time event \\
   Down Time \#2 Operator Set ID & String & 15 & Operator needed from
   this set \\
    & & & during this down time event\\
   Down Time \#3 Description & String & 25 & Description of tool down time
   type 3 \\
   Down Time \#3 Type & Integer & 5 & 0 = time based; 1 = run based \\
   Time or Runs Between \#3 & Float & 15 & Mean time (runs) between this \\
    & & & down time event \\
   Duration \#3 & Float & 10 & Duration of this down time event \\
   Down Time \#3 Operator Set ID & String & 15 & Operator needed from
   this set \\
    & & & during this down time event\\
   Down Time \#4 Description & String & 25 & Description of tool down time
   type 4 \\
   Down Time \#4 Type & Integer & 5 & 0 = time based; 1 = run based \\
   Time or Runs Between \#4 & Float & 15 & Mean time (runs) between events \\
   Duration \#4 & Float & 10 & Duration of this down time event \\
   Down Time \#4 Operator Set ID & String & 15 & Operator needed from
   this set \\
    & & & during this down time event\\
   Down Time \#5 Description & String & 25 & Description of tool down time
   type 5 \\
   Down Time \#5 Type & Integer & 5 & 0 = time based; 1 = run based \\
   Time or Runs Between \#5 & Float & 15 & Mean time (runs) between events \\
   Duration \#5 & Float & 10 & Duration of this down time event \\
   Down Time \#5 Operator Set ID & String & 15 & Operator needed from
   this set \\
    & & & during this down time event\\ \hline
   \end{tabular}
   \end{center}
   
   \newpage
   \subsection{Operator Set File}
   The format for operator set records is similar to the format for
   the tool set.
   
   \begin{center}
   \begin{tabular}{|l|c|c|l|}  \hline
   {\em Field Name} &  {\em Type} & {\em Width } & {\em Remarks} \\ \hline
   Operator Set ID & String & 15 & Operator set identifier \\
   Operator Description & String & 25 & Operator Set Name \\
   Quantity  & Integer & 5 & Number of operators in this set \\
   Break \#1 Description & String & 25 & Description of break type \\
   Time Between \#1 & Float & 15 & Time between breaks of this type \\
   Duration \#1 & Float & 10 & Duration of breaks of this type \\
   Break \#2 Description & String & 25 & Description of break type \\
   Time Between \#2 & Float & 15 & Time between breaks of this type \\
   Duration \#2 & Float & 10 & Duration of breaks of this type \\
   Break \#3 Description & String & 25 & Description of break type \\
   Time Between \#3 & Float & 15 & Time between breaks of this type \\
   Duration \#3 & Float & 10 & Duration of breaks of this type \\ \hline
   \end{tabular}
   \end{center}
   
   \subsection{Volume/Release  File}
   
   \begin{center}
   \begin{tabular}{|l|c|c|l|}  \hline
   {\em Field Name} &  {\em Type} & {\em Width } & {\em Remarks} \\ \hline
   Process Flow & String & 10 & Process Flow ID \\
   Product ID & String & 10 & Unique product family ID  \\
   Product Name & String & 25 & Name of Product (optional) \\
   Start Rate & Float & 10 & Number of wafers per day  released into \\
    & & & line based on 7 day/week operation \\
   Lot Size & Integer & 10 & Number of wafers in a released lot \\ \hline
   \end{tabular}
   \end{center}
  
  \subsection{Comments File}
   The comments file contains information about the factory represented 
   by the dataset, including the type of product manufactured, 
   number of process flows
   in the dataset, number of equipment groups, and approximate wafer starts
   per month. The file also contains sample simulation output, including  
   assumptions, simulation run conditions, and results (cycle times and 
   number of outs).




 
   \newpage
   \section{Additional Information}
   
   In this section, we list, in no particular order,
   relevant issues and clarifying points
   relating to the data sets.
   
   \begin{enumerate}
   \item All times are specified in minutes except if stated explicitly
   otherwise.
   \item  All tools in a tool set and all operators in a operator set
   are considered identical.  In particular, all tools in a tool set
   are considered qualified to perform all operations coming to that
   tool set.  The same holds true for operators.
   \item  An operation is uniquely specified by its Process ID
   (or Rework ID) and Step ID.
   However, we allow multiple entries
   for the same Product ID, Step ID pair.  The resolution of such
   multiple entries should be made in the description field.
   Multiple entries might be used to specify alternative tool sets
   or operator sets that can be used to perform an operation.
   \item The following formulas can be used to calculate processing
   time per lot (pt), time until tool becomes free (
   tf), and
   total lot cycle time through an operation (ct): \\
   pt = \verb+Time per Batch+ * Number of batches required for the lot \\
   ~~~~~+ \verb+Time per Lot+  +  \verb+Time per Wafer in Process+
    * Number of wafers in lot \\
   ~~~~~+ \verb+Product Setup+ (if appropriate) +
   \verb+Group Setup+ (if appropriate) \\
   tf = \verb+Load Time+ + pt + \verb+Unload Time+ \\
   ct = \verb+Load Time+ + pt + \verb+Wafer Travel Time+ +
   \verb+Unload Time+ \\
   
   \item No distributional information is included in the data
   sets beyond first moment (mean) information.  Fields with
   names like \verb+Load Time+ have the implicit prefix qualifier
   \verb+Mean+.
   \item  The wafer start rates given in the Volume/release
   file are intended as guidelines.  The exact method by
   which lots are released is to be determined by the users of the
   data.
   \item  Information on
   process holds, engineering holds and send aheads is not included.
   \item Time bound sequences cannot be specified explicitly.
   \item The following types of tools typically found in a semiconductor
   wafer manufacturing line can be modeled by making appropriate
   use of the processing time parameters:
   \begin{enumerate}
   \item Single wafer tools.  Set \verb+Time per Wafer in Process+ field
   to wafer processing time and all other process time parameters to 0.
   \item Batch tools.  Set \verb+Min Batch Size+ and \verb+Max Batch Size+ 
   fields appropriately.  Set \verb+Time per Batch+
   field to batch processing time and all other process time
   parameters to 0. Lots with the same batch ID may be batched 
   together (up to the max batch size). Where batch IDs are left blank
   only lots at the same step of the same process flow may be batched together.
   \item Multi-Sequence tools.  Set \verb+Min Batch Size+ and 
   \verb+Max Batch Size+ fields appropriately.   Set \verb+Time per Batch+
   field to the largest single tank time and set 
   \verb+ Wafer Travel Time+ field to the  remaining tank time. Set other process
   time parameters to 0.
   \item Conveyor tools. Set \verb+Time per Wafer in Process+ field
   to wafer processing time and \verb+ Wafer Travel Time+ field 
   to the wafer travel time.  Set other process time parameters to 0.
   \item Inspect tools. Set \verb+Time per Lot+ field to inspect time.
   Set other process time parameters to 0.
   \item Certain cluster tools.  Model as batch tools.
   \item Linked track/ linked lithography tools. Model as conveyor
   tools.
   \end{enumerate}
   \item Group and spec. setups are assumed to be done whenever the group or 
   process flow step ID changes from one operation to another. However, the duration of the 
   setup does not vary as a function of the specific group or product preceding 
   the operation.
   \item An Operator Set ID or Tool Set ID field that is blank indictes that no resource is 
   needed for that step.
   \end{enumerate}
   \newpage
   
   \end{document}
